\scanpage{15}
\mspage{13}

\uline{Dém} a) Comme $\pi'_0\longrightarrow\pi_0$ est évidemment un isom (car $k'/k$ est \unsure{l'ext. triviale}) il suffit de prouver que $I'\simeq I$, et comme $I'_t\xrightarrow{\sim}I_t$ (car $e=1$) il suffit de prouver que $P'\simeq P$.

b) On sait que $\pi'\mathrel{\text{---}}\pi$ est surjectif, i.e pour tte ext. finie sép. $K_1/K$, $K_1\otimes_K\widehat K$ est un \uline{corps} : en effet, la clôture normale $V_1$ de $V$ dans $K_1$ est finie, donc $V_1\otimes_V\widehat V\simeq\widehat V_1$, \deleted{est} \inserted{c'est donc} un anneau de val. discrète ($V_1$ en étant un), donc $K_1\otimes_K\widehat K\simeq\widehat V_1\otimes_V\widehat K$ est un corps.

\deleted{\unsure{D'ailleurs}} \deleted{Quitte à remplacer $V,V'$ par $\widetilde V,\widetilde V'$ et à comparer leurs complétés \unsure{$\widehat{\widetilde V}\simeq\widehat{\widetilde V'}$} de ces deux anneaux, \unsure{on est ramené}}

c) On en conclut, sur les \inserted{hom de} suites exactes, que $I'\mathrel{\text{---}}I$ est surjectif, puis $P'\mathrel{\text{---}}P$ \uline{surjectif}.
Il reste à prouver que $P'\mathrel{\text{---}}P$ est injectif, et pour ceci que $\pi'\mathrel{\text{---}}\pi$ est injectif, \deleted{\unsure{On}} i.e. que toute extension séparable finie $K'_1$ de $K'=\widehat K$ provient d'une ext. sép. finie de $K$. Or ceci résulte du fait connu (``continuité des racines p. r. aux coefficients'') que deux équations séparables \unsure{\deleted{de $K$}|\inserted{de $K$}} à coefficients dans $\widehat K$ \uline{voisins} définissent des \deleted{\unsure{corps}} $\widehat K$-algèbres isomorphes.
\margin{14}
